% Abstract. No number that is not in the body. Every number has a % src.
\begin{abstract}
AlphaFold2 is written in JAX. JAX treats each chip of a tensor processing unit (TPU) slice as a separate device. A JAX
program does not spread its work over the chips by itself, so a user who runs AlphaFold2 on a slice can pay for chips
that do nothing. We found no published measurement of how AlphaFold2 inference uses a multi-chip TPU slice. Our goal
was to measure how much of an eight-chip TPU v5e slice is occupied by AlphaFold2 inference and to find which way of
mapping the work in JAX uses every chip.
All 12 experiments ran \texttt{model\_3} with random weights, a single-sequence multiple sequence alignment, and no recycling, so we studied execution and not prediction quality. % src: recomputed: COUNT.tpu_experiments_with_raw_data; src/spike_tpu_forward_pass.py ("model_3")
By default, only one chip held data. We call each input sequence a query.
Batching up to eight queries of 118 residues on one chip with \texttt{jax.vmap} did not raise throughput. % src: results/sweep/batching_sweep.json -> [*].batch_size; recomputed: INPUT.toy_length
With \texttt{jax.pmap}, the eight chips ran eight independent queries in parallel.
On a grid of 1 to 8 chips and 100 to 1000 residues, the full slice gave 6.5 to 7.9 times the throughput of one chip. % src: results/sweep/scaling_law_data.json -> grid[*].chips, grid[*].length; recomputed: V2.speedup_8over1_min, V2.speedup_8over1_max
Without sharding annotations, JAX's automatic partitioning did not make one query run faster. Most settings ran once,
and the logs of the TPU runs were not kept.
\end{abstract}

\begin{IEEEkeywords}
AI accelerators, AlphaFold2, bioinformatics, JAX, parallel processing, performance evaluation, scalability
\end{IEEEkeywords}
